% ==============================================================================
% capitulos/06-analisis-diseno.tex - Análisis y Diseño del Sistema
% ==============================================================================

\chapter{Análisis y Diseño del Sistema}
\label{cap:analisis_diseno}

\section{Modelo del Alcance}
El proyecto comprende el desarrollo de una solución integral para la Defensoría de los Derechos Politécnicos. El alcance se define mediante los siguientes límites operativos:

\begin{itemize}
    \item \textbf{Funcional:} Cubre desde la recepción de la queja (Primer Contacto) hasta la emisión del dictamen final. Incluye la clasificación asistida por IA y la gestión documental digital.
    \item \textbf{Poblacional:} Exclusivo para la comunidad del IPN (alumnos, docentes y administrativos) con credenciales institucionales.
    \item \textbf{Tecnológico:} Implementación de una arquitectura de microservicios desplegada mediante contenedores para asegurar la disponibilidad.
\end{itemize}

\section{Modelo del Negocio}
El modelo de negocio se basa en la digitalización del Manual de Procedimientos de la DDP. El flujo principal se describe a continuación:



\begin{enumerate}
    \item \textbf{Captación:} El usuario ingresa la queja; el sistema valida la identidad vía correo institucional.
    \item \textbf{Pre-clasificación:} El motor de PLN analiza el texto y sugiere si el caso es "Competente" o "No Competente".
    \item \textbf{Asignación:} Un abogado asesor valida la clasificación y radica el expediente.
    \item \textbf{Seguimiento:} El sistema gestiona las evidencias y notifica cambios de estado al quejoso mediante folios digitales.
\end{enumerate}

\section{Modelo Dinámico}
Para representar el comportamiento del sistema ante eventos externos, se definen los diagramas de secuencia para los casos de uso críticos.

\subsection{Diagrama de Secuencia: Registro y Clasificación}
Este modelo describe la interacción entre el Frontend (Angular), el API Gateway (Nginx), el Microservicio de Gestión (Spring Boot) y el Microservicio de IA (Python).



\section{Análisis de Riesgo}
Se identifican los riesgos que podrían afectar la viabilidad del proyecto utilizando una matriz de impacto y probabilidad.

\begin{table}[htbp]
\centering
\small
\begin{tabular}{|p{3cm}|c|c|p{6cm}|}
\hline
\textbf{Riesgo} & \textbf{Prob.} & \textbf{Impacto} & \textbf{Estrategia de Mitigación} \\ \hline
Baja precisión del modelo de IA & Media & Alto & Implementar un modo "asistido" donde el abogado siempre valida el resultado. \\ \hline
Vulneración de datos sensibles & Baja & Crítico & Cifrado de base de datos y uso de protocolos HTTPS/SSL gestionados en Nginx. \\ \hline
Curva de aprendizaje de microservicios & Media & Medio & Uso de Docker para estandarizar entornos y documentación exhaustiva. \\ \hline
\end{tabular}
\caption{Matriz de Riesgos del Proyecto.}
\end{table}

\section{Costos del Proyecto}
Dada la naturaleza del Trabajo Terminal, los costos se dividen en infraestructura técnica y capital humano proyectado.

\begin{itemize}
    \item \textbf{Infraestructura:} Servidor institucional o instancia Cloud (aprox. \$200 - \$500 USD anuales).
    \item \textbf{Software:} \$0.00 (Uso de tecnologías de código abierto: Java, Python, PostgreSQL, Linux).
    \item \textbf{Capital Humano:} Estimación de 1,200 horas de desarrollo/investigación por parte de los tres integrantes.
\end{itemize}


